\documentclass[11pt]{article}

\usepackage[margin=1in]{geometry}
\usepackage{amsmath,amssymb}
\usepackage{mathtools}
\usepackage{microtype}

\title{Root-Free Fractions Whose Maxima Have a Square-Root Denominator}
\author{}
\date{}

\begin{document}

\maketitle

\section{Goal}

We want to identify a root-free rational function in \(p\) and \(a\) whose maximum over \(p\) can be written in the form
\[
1+\frac{c}{da+\sqrt{a(ea+f)}}.
\]
Here \(a\) is the variable parameter, while \(p\) is the variable being optimized over.

\section{The root-free fraction}

A natural class of fractions with this property is
\[
\boxed{
G(p,a)
=
\frac{\alpha a p^2+Bp+\beta a+\gamma-q}
{\alpha a p^2+\beta a+\gamma}.
}
\]
Here
\[
\alpha,\ B,\ \beta,\ \gamma,\ q
\]
are constants.

This fraction has no square roots. Its structure is:
\begin{itemize}
    \item the numerator and denominator have the same quadratic term \(\alpha a p^2\);
    \item the numerator has a linear term \(Bp\);
    \item the denominator has affine constant part \(\beta a+\gamma\);
    \item the numerator constant part differs from the denominator constant part by the fixed amount \(q\).
\end{itemize}

Equivalently,
\[
G(p,a)
=
\frac{A(a)p^2+Bp+C(a)}
{A(a)p^2+D(a)},
\]
where
\[
A(a)=\alpha a,
\qquad
D(a)=\beta a+\gamma,
\qquad
C(a)=D(a)-q.
\]

\section{Maximization over \(p\)}

For the general form
\[
G(p)
=
\frac{Ap^2+Bp+C}{Ap^2+D},
\]
the maximizing critical point is
\[
p_\ast
=
\frac{D-C+\sqrt{(D-C)^2+\frac{B^2D}{A}}}{B}.
\]
At that critical point,
\[
G_{\max}
=
1+\frac{B}{2Ap_\ast}.
\]
Substituting the expression for \(p_\ast\) gives
\[
G_{\max}
=
1+
\frac{B^2}
{2A\left(D-C+\sqrt{(D-C)^2+\frac{B^2D}{A}}\right)}.
\]

Now insert
\[
A=\alpha a,
\qquad
D=\beta a+\gamma,
\qquad
D-C=q.
\]
Then
\[
G_{\max}(a)
=
1+
\frac{B^2}
{2\alpha a\left(q+\sqrt{q^2+\frac{B^2(\beta a+\gamma)}{\alpha a}}\right)}.
\]
Multiplying the outer \(a\) into the square-root term gives
\[
\boxed{
G_{\max}(a)
=
1+
\frac{B^2}
{
2\alpha q\,a
+
2\sqrt{\alpha a\left((\alpha q^2+B^2\beta)a+B^2\gamma\right)}
}.
}
\]

\section{Writing it in the desired structure}

The previous expression is already of the form
\[
1+\frac{c}{da+\sqrt{a(ea+f)}}.
\]
Indeed,
\[
G_{\max}(a)
=
1+
\frac{B^2}
{
2\alpha q\,a
+
\sqrt{
a\left(
4\alpha(\alpha q^2+B^2\beta)a
+
4\alpha B^2\gamma
\right)
}
}.
\]
Therefore the parameter correspondence is
\[
\boxed{
c=B^2,
}
\]
\[
\boxed{
d=2\alpha q,
}
\]
\[
\boxed{
e=4\alpha(\alpha q^2+B^2\beta),
}
\]
and
\[
\boxed{
f=4\alpha B^2\gamma.
}
\]

So the root-free fraction
\[
\boxed{
\frac{\alpha a p^2+Bp+\beta a+\gamma-q}
{\alpha a p^2+\beta a+\gamma}
}
\]
has maximum
\[
\boxed{
1+\frac{c}{da+\sqrt{a(ea+f)}}
}
\]
with
\[
\boxed{
c=B^2,\qquad
d=2\alpha q,\qquad
e=4\alpha(\alpha q^2+B^2\beta),\qquad
f=4\alpha B^2\gamma.
}
\]

\section{Example}

The earlier expression reduced to
\[
F(p,a)
=
\frac{
4ap^2+\frac{12}{5}p+\frac{8+74a}{25}
}{
4ap^2+\frac{38+74a}{25}
}.
\]
This matches the template with
\[
\alpha=4,
\qquad
B=\frac{12}{5},
\qquad
\beta=\frac{74}{25},
\qquad
\gamma=\frac{38}{25},
\qquad
q=\frac65.
\]
Hence
\[
c=B^2=\frac{144}{25},
\]
\[
d=2\alpha q
=2\cdot4\cdot\frac65
=\frac{48}{5},
\]
\[
e=4\alpha(\alpha q^2+B^2\beta)
=4\cdot4\left(4\cdot\frac{36}{25}+\frac{144}{25}\cdot\frac{74}{25}\right),
\]
and
\[
f=4\alpha B^2\gamma
=4\cdot4\cdot\frac{144}{25}\cdot\frac{38}{25}.
\]
These constants can be divided by a common scaling factor inside the denominator.
After simplifying the overall fraction, one obtains the cleaner equivalent form
\[
\boxed{
F_{\max}(a)
=
1+\frac{3}{5a+\sqrt{a(99a+38)}}.
}
\]

\section{Summary}

The important root-free template is
\[
\boxed{
G(p,a)
=
\frac{\alpha a p^2+Bp+\beta a+\gamma-q}
{\alpha a p^2+\beta a+\gamma}.
}
\]
It produces the square-root maximum structure because:
\begin{enumerate}
    \item optimizing over \(p\) leads to a quadratic equation;
    \item solving that quadratic introduces a square root;
    \item the coefficient of \(p^2\) is proportional to \(a\);
    \item the denominator's constant part is affine in \(a\);
    \item the numerator and denominator differ by a constant \(q\) in their non-\(p\) parts.
\end{enumerate}

\end{document}
